\documentclass[11pt]{article}
\usepackage[utf-8]{inputenc}
\usepackage{amsmath}
\usepackage{amssymb}
\usepackage{booktabs}
\usepackage{xcolor}
\usepackage{hyperref}

\title{Chaboche-v1 Level-2 Cycle-Jump Preparation Summary}
\author{Thesis Work}
\date{}

\begin{document}

\maketitle

\section{Overview}

This document synthesizes the progression from Level-1 scalar SDV1 cycle-jump postprocessing to Level-2 preparation diagnostics for the simplified Chaboche-v1 unified viscoplastic UMAT. It explains the current findings and the rationale for deferring full STATEV injection to a later phase.

\section{STATEV Inventory}

The Chaboche-v1 UMAT manages 15 solution-dependent state variables (STATEV):

\begin{table}[ht]
\centering
\small
\begin{tabular}{r|l|l|l}
\toprule
\textbf{Index} & \textbf{Symbol} & \textbf{Physical meaning} & \textbf{Category} \\
\midrule
1  & $p$ & Accumulated viscoplastic strain & Active / required \\
2--4 & $X_{11}, X_{22}, X_{33}$ & Backstress normal components & Active / required \\
5--7 & $X_{12}, X_{13}, X_{23}$ & Backstress shear components & Near-zero \\
8--10 & $E_{vp,11}, E_{vp,22}, E_{vp,33}$ & Visc.~strain normal & Active / required \\
11--13 & $E_{vp,12}, E_{vp,13}, E_{vp,23}$ & Visc.~strain shear & Near-zero \\
14 & RISO & Isotropic hardening stress & Recomputable \\
15 & DP & Last visc.~multiplier increment & Diagnostic \\
\bottomrule
\end{tabular}
\caption{STATEV inventory for Chaboche-v1 UMAT. STATEV(1--13) form the independent material state.}
\end{table}

\section{Level-1: Scalar SDV1 Postprocessing Prediction}

Extract the accumulated viscoplastic strain $\text{STATEV}(1)$ at each cycle end.  Use cycles 2--10 as the stabilized reference window. Estimate the per-cycle increment and extrapolate.

\subsection{Validation Result}
\begin{itemize}
  \item Predicted $\text{SDV1}$ at cycle 20 (from cycle 10 base): $0.142121435146$
  \item Explicit Abaqus $\text{SDV1}$ at cycle 20: $0.1420256943$
  \item Relative error: $0.0674\%$
\end{itemize}

Conclusion: scalar SDV1 cycle jumping is accurate and thesis-ready.

\section{Level-2a: Full-State Cycle-History Extraction}

Extract entire state vector $\text{STATEV}(1\text{--}15)$ at each cycle end. Perform stability classification.

\subsection{Key Findings}
\begin{itemize}
  \item $\text{STATEV}(1)$: \textbf{stable extrapolation candidate}
  \item $\text{STATEV}(2\text{--}4)$: \textbf{needs caution} (high cycle-to-cycle variation)
  \item $\text{STATEV}(5\text{--}7)$: \textbf{near-zero}
  \item $\text{STATEV}(8\text{--}10)$: \textbf{needs caution}
  \item $\text{STATEV}(11\text{--}13)$: \textbf{near-zero}
  \item $\text{STATEV}(14\text{--}15)$: \textbf{recomputable/diagnostic}
\end{itemize}

\section{Level-2b: Vector-Valued STATEV Cycle-Jump Control}

Extend scalar cycle-jump to vector: $\text{STATEV}(1), \text{STATEV}(2\text{--}4), \text{STATEV}(8\text{--}10)$. Use adaptive curvature checking for conservative global jump.

\subsection{Original Result}
\begin{itemize}
  \item Conservative global $\Delta N$: \textbf{2}
  \item Adaptive target cycle: \textbf{12}
  \item Controlling component: \textbf{STATEV(2) / $X_{11}$}
  \item First-order SDV1 error at target: $0.0118\%$
\end{itemize}

\section{Level-2c: Exact-Output Diagnostic Branch}

Request exact integer cycle-end frames using Abaqus \texttt{TIME MARKS=YES}. Compare against original nearest-frame run.

\subsection{Exact-Output Results}
\begin{itemize}
  \item Max cycle-end time error: $0$ (exact frames) vs. original $0.00974$
  \item Cycle-20 $\text{STATEV}(1)$: $0.134750679135$ (exact) vs. $0.142025694251$ (original)
  \item Relative difference: \textbf{5.12\%}
\end{itemize}

\textit{Key insight}: \texttt{TIME MARKS=YES} forced smaller time increments, altering the UMAT-integrated response. The Chaboche-v1 UMAT is \textbf{increment-schedule sensitive}.

\subsection{Exact Vector Result}
\begin{itemize}
  \item Conservative global $\Delta N$: \textbf{1} (vs. original 2)
  \item Adaptive target: \textbf{11} (vs. original 12)
  \item Controlling component: \textbf{STATEV(2) / $X_{11}$} (unchanged)
\end{itemize}

\section{Final Decision: Deferring STATEV Injection}

\subsection{Status}
\begin{itemize}
  \item Level 1: validated ($0.0674\%$ error); thesis-ready
  \item Level 2a: completed; full-state stability known
  \item Level 2b: completed; vector approach more conservative
  \item Level 2c: completed; increment-schedule sensitivity revealed
\end{itemize}

\subsection{Why Not Proceed to Injection}

\begin{enumerate}
  \item \textbf{Increment-schedule sensitivity}: Injected state at a given cycle may produce different predictions depending on next Abaqus run's integration scheme.
  \item \textbf{Vector-valued complexity}: Consistent treatment of 13 components requires robust integration; current UMAT is not robust enough.
  \item \textbf{Thesis narrative}: Level-2 diagnostics are strong content; stopping here allows clear discussion of why robustness is necessary.
\end{enumerate}

\section{Distinction Between Method Levels}

\begin{description}
  \item[Level 1:] Postprocessing prediction from existing ODB. Read-only. No UMAT changes. Accuracy: postprocessing-level.
  \item[Level 2:] Full-state extraction, stability analysis, increment-schedule diagnostics. Test runs only. No UMAT changes. Accuracy: diagnostic.
  \item[Level 3:] Abaqus restart/state-injection with STATEV modification. Full forward integration. Subject to UMAT robustness. \textit{Not pursued in this project.}
\end{description}

\section{Thesis Narrative}

The Chaboche-v1 unified viscoplastic UMAT cycle-jump workflow demonstrates:

\begin{enumerate}
  \item Scalar postprocessing cycle jumping is accurate ($0.0674\%$ error).
  \item Vector-valued jumping is feasible but more conservative, requiring robust integration.
  \item Integration robustness is a prerequisite for restart injection.
  \item Exact phase-point output can remove frame-time ambiguity but at the cost of altering the increment schedule.
\end{enumerate}

By stopping at Level-2 preparation, the thesis presents a rigorous diagnostics and preparation study without committing to advanced continuation that could amplify robustness issues.

\end{document}
